\begin{abstract}
The single-source capacitated facility location problem (SSCFLP) asks which facilities to
open and which single facility serves each customer, under capacity limits. Even instances
with 30 candidate facilities and 150 customers are not solved to optimality by a modern
open-source MILP solver within a minute, because the linear relaxation leaves a gap of
1.6--4.4\% when fixed costs dominate. Machine learning has been proposed to reduce the search
space of such problems, by pruning variables, ranking them, or choosing which part of a solution
to re-optimize. We ask how much of the reported benefit is due to learning and how much to the
underlying search-space reduction, by pairing every learned component with a learning-free
counterpart of the same function under equal wall-clock budgets.

We study three mechanisms. (i) \emph{Pruning} by a facility ranking is unsafe under capacity
coupling: keeping the best known solution reachable in 80\% of the instances required retaining
about 80\% of the facilities, for a graph neural network (GNN) and for the linear relaxation
alike. (ii) An \emph{adaptive expansion} solves nested restricted problems over the ranked
facilities and stops when reduced costs certify that no excluded facility can improve the
incumbent, which yields a proven optimum without solving the full model. (iii) A
\emph{cluster-oriented large neighborhood search} re-optimizes one group of customers at a time
with the rest fixed and capacities discounted; we prove that every accepted move is globally
feasible and charges each fixed cost once, whereas independently solved clusters violated
capacity in every examined instance. A transposition-style memory identifies subproblems that
have already failed.

In validation pilots, all restriction-based methods had a primal integral two to four times lower
than the solver on the full model, and the hybrid of (ii) and (iii) finished 0.5--0.9\% from the
best known solution against 2.4--2.7\%. Learning, however, was not the source of the gain:
ranking facilities by the linear relaxation was numerically at least as good as the GNN in every
method (differences not significant on 16 instances), and
neither a learned nor a GNN-guided subproblem selector separated from a plain rotation. More
than half of the subproblems selected by the neighborhood search were repetitions of already
failed ones. \todo{one sentence with the closed-test result and its significance}. The
experiments use pinned cores, two clocks, frozen data and pre-registered hypotheses, and we also
report what did not transfer when we reconstructed five recent learning-for-optimization
approaches in an open-source stack.
\end{abstract}
